\chapter{Struttura della Lega}
\label{cap:struttura-lega}

\section{Finalità e principi ispiratori}
\label{art:1.1}

La Lega FantaScarrupat è un'associazione amatoriale a carattere ludico, ispirata ai valori dell'amicizia, del rispetto reciproco e della partecipazione condivisa. La Lega si fonda su principi democratici: ogni squadra partecipa con pari dignità alla vita della competizione, contribuendo attivamente alle decisioni comuni.\\
La gestione del gioco è affidata alla piattaforma digitale “Fantacalcio.it”.

\section{Composizione e ruolo del Consiglio di Lega}
\label{art:1.2}

Tutti i fantallenatori iscritti alla Lega fanno parte del Consiglio di Lega, organo deliberativo collegiale incaricato di adottare le decisioni strategiche, organizzative e regolamentari. Il Consiglio stabilisce le regole della stagione, convoca e gestisce le aste, e risolve eventuali controversie.

Per la stagione sportiva 2026/2027, il Consiglio risulta composto da otto membri. La Presidenza è affidata ad Antonio Di Giorgio, che ricopre altresì la funzione di Amministratore di Lega. Ciascun membro del Consiglio dispone di pari diritti di voto e doveri di partecipazione.

\section{Il Presidente della Lega}
\label{art:1.3}

Il Presidente è il garante del corretto svolgimento delle attività della Lega. Oltre a convocare e presiedere le riunioni, egli ha facoltà di:
\begin{itemize}
\item indire le votazioni e sottoporre proposte al Consiglio;
\item segnalare comportamenti contrari allo spirito sportivo della Lega;
\item proporre la sospensione temporanea di un membro dalle votazioni, qualora quest'ultimo assuma atteggiamenti offensivi, ostruzionistici o in aperta violazione del regolamento, previa comunicazione motivata al Consiglio, secondo la procedura di cui all'articolo \ref{art:2.5}.
\end{itemize}

\section{Il Tesoriere}
\label{art:1.4}

La funzione di Tesoriere è attribuita dal Consiglio di Lega e, per la stagione sportiva 2026/2027, è ricoperta da Antonio Di Giorgio, che la cumula con la carica di Presidente.\\
Il Tesoriere custodisce le somme versate a titolo di quota, rilascia le ricevute nominative di cui all'articolo \ref{art:10.1}, amministra il fondo di riserva di cui all'articolo \ref{art:10.5} e presenta al Consiglio, al termine di ciascuna stagione, un rendiconto delle entrate e delle uscite.

\section{Ritiro o esclusione di un partecipante}
\label{art:1.5}

Il numero delle squadre partecipanti resta invariato per l'intera stagione, poiché su di esso sono generati i calendari e i tabelloni di tutte le competizioni (vedi Titolo \ref{cap:competizioni}).

Qualora un partecipante si ritiri o venga escluso a stagione iniziata, la sua squadra non viene rimossa e prosegue la stagione in una delle seguenti forme, scelta dal Consiglio:

\begin{itemize}
\item \textbf{squadra affidata}: un altro fantallenatore ne assume la gestione in aggiunta alla propria e ne versa la quota; la squadra concorre a pieno titolo a tutte le competizioni, premi compresi;
\item \textbf{squadra cuscinetto}: nessuno ne assume la gestione e lo schieramento è curato dall'Amministratore di Lega secondo le modalità indicate più avanti; la squadra concorre alle competizioni ai soli fini dei risultati delle altre squadre e non ha diritto ai premi, che restano acquisiti al fondo di riserva (vedi \ref{art:10.5}). La relativa quota, se non versata dal partecipante uscente, è ripartita tra gli altri membri ai sensi dell'articolo \ref{art:3.3}, affinché il montepremi resti invariato.
\end{itemize}

In entrambi i casi la squadra conserva crediti e risultati già maturati; la rosa è del pari conservata, salvo che il Consiglio opti per la modalità 3 della disciplina che segue.

Il fantallenatore che gestisce due squadre non può effettuare scambi tra di esse, dispone comunque di un solo voto in Consiglio e, in caso di incontro diretto tra le proprie squadre, comunica entrambe le formazioni nel canale ufficiale della Lega prima dell'inizio della giornata.

Quanto alla squadra cuscinetto, il Consiglio sceglie una delle seguenti modalità di gestione:

\begin{enumerate}
\item \textbf{Rosa invariata, schieramento fisso}: la rosa resta quella del partecipante uscente e l'Amministratore di Lega compone una sola volta formazione e panchina, scegliendo i calciatori meno impiegati nel campionato reale. Lo schieramento resta immutato per tutte le giornate residue della stagione.
\item \textbf{Rosa invariata, schieramento per giornata}: la rosa resta quella del partecipante uscente e l'Amministratore di Lega compone lo schieramento giornata per giornata, secondo il medesimo criterio. Nella giornata in cui la squadra dell'Amministratore affronta la cuscinetto, lo schieramento è composto da altro membro designato dal Consiglio.
\item \textbf{Rosa sostituita}: l'Amministratore di Lega ricompone la rosa, nel rispetto della consistenza prevista dall'articolo \ref{art:4.2} (30 calciatori, di cui almeno 3 portieri), con calciatori presenti nel listone ufficiale di Serie A privi di impiego, così da evitare che la squadra raggiunga la soglia di cui all'articolo \ref{art:5.2}, e ne fissa una sola volta lo schieramento. I calciatori liberati per effetto della sostituzione non rientrano nella lista svincolati e non sono acquisibili dalle altre squadre per il resto della stagione.
\end{enumerate}

In ogni modalità lo schieramento è comunicato nel canale ufficiale della Lega prima dell'inizio della giornata cui si riferisce. Alla squadra cuscinetto non si applicano i fattori e le funzioni opzionali di cui al Titolo \ref{cap:modalita-calcolo}; essa non partecipa alle aste di riparazione, non effettua scambi né prestiti (vedi \ref{art:6.5}) e non può avvalersi delle Priorità (vedi \ref{art:6.4}) né delle sostituzioni per infortunio (vedi \ref{art:7.3}). Per il resto valgono le ordinarie regole di sostituzione di cui all'articolo \ref{art:4.5}.

Ne consegue che la squadra cuscinetto difficilmente raggiunge la soglia di 66 punti (vedi \ref{art:5.2}). L'incontro con essa non costituisce pertanto una vittoria automatica: chi la affronta prevale soltanto realizzando a propria volta almeno un gol, e in difetto l'incontro termina in parità.

\section{Sorteggi e determinazione degli ordini}
\label{art:1.6}

La Lega dispone del software «Ruota della Fortuna», accessibile tramite il sito ufficiale (vedi \ref{art:11.1}) e avviato dall'Amministratore di Lega.

Il suo utilizzo non è obbligatorio: i partecipanti possono sempre concordare direttamente fra loro l'ordine o la scelta da operare. Vi si ricorre quando un accordo non venga raggiunto, ovvero quando si preferisca affidarsi fin da subito a un criterio casuale.

Essa costituisce il metodo di riferimento ogniqualvolta il presente regolamento richieda di stabilire un ordine, di operare una scelta casuale o di dirimere una situazione di parità non altrimenti disciplinata, e le parti non abbiano convenuto diversamente. Rientrano fra tali ipotesi:

\begin{itemize}
\item l'ordine di chiamata dei fantallenatori nelle aste a chiamata diretta libera (vedi \ref{art:4.3.2});
\item l'ordine di intervento nella fase integrativa di chiamata diretta (vedi \ref{art:4.3.4}) e nelle aste di riparazione (vedi \ref{art:6.1});
\item l'individuazione del fantallenatore cui affidare una squadra ai sensi dell'articolo \ref{art:1.5}, qualora più di uno si renda disponibile e il Consiglio non deliberi diversamente;
\item ogni altra ipotesi in cui occorra un sorteggio e il regolamento non preveda espressamente un metodo diverso.
\end{itemize}

L'estrazione è eseguita dall'Amministratore di Lega alla presenza di almeno un altro membro del Consiglio, ovvero in diretta nel canale ufficiale della Lega, e il suo esito è comunicato a tutti i partecipanti. L'esito del sorteggio è definitivo; eventuali contestazioni sulla regolarità della procedura sono proposte con le modalità di cui all'articolo \ref{art:2.7}.

Resta fermo quanto previsto per la generazione dei calendari e dei tabelloni delle competizioni, disciplinata dal Titolo \ref{cap:competizioni}.


\chapter{Votazioni di Lega}
\label{cap:Votazioni}

\section{Disposizioni comuni}
\label{art:2.1}

Ai fini del presente regolamento si intendono per \textbf{membri aventi diritto} i membri del Consiglio che prendono parte alla singola votazione, con esclusione tanto di coloro che ne siano esclusi ai sensi dell'articolo \ref{art:2.5}, quanto di coloro che non abbiano risposto entro il termine di cui all'articolo \ref{art:2.3}. Ogni maggioranza prevista dal presente regolamento si calcola esclusivamente su tale base, che è pertanto variabile in funzione della singola deliberazione. I risultati frazionari si arrotondano all'unità superiore.

Il voto è palese. Nelle deliberazioni riguardanti le persone, e in particolare nelle sanzioni, su richiesta anche di un solo membro si procede a scrutinio segreto; in tal caso si rende noto il solo conteggio.

L'astensione espressa entro il termine costituisce partecipazione alla votazione: il membro che si astiene è computato tra i membri aventi diritto ma non tra i voti espressi.

Non è richiesta alcuna verbalizzazione: l'esito della votazione è comunicato nel canale ufficiale della Lega.

\section{Maggioranza e parità di voti}
\label{art:2.2}

Tutte le deliberazioni del Consiglio, qualunque ne sia l'oggetto, sono approvate con la \textbf{maggioranza assoluta dei membri aventi diritto}, ossia con un numero di voti favorevoli superiore alla metà di essi. Con otto aventi diritto occorrono pertanto cinque voti favorevoli. Non esistono materie soggette a maggioranze diverse: la soglia è unica per l'intero regolamento.

In caso di parità si distinguono due ipotesi.

Se la votazione ha ad oggetto la modifica, la revoca o la sostituzione di una disciplina già in vigore, la proposta si intende respinta e resta in vigore la disciplina precedente.

Se invece la votazione ha ad oggetto una scelta nuova, ovvero una questione che debba necessariamente avere un esito e per la quale non esista una disciplina precedente cui tornare, la parità è risolta mediante la «Ruota della Fortuna» (vedi \ref{art:1.6}) tra le sole opzioni che hanno riportato il maggior numero di voti.

\section{Svolgimento delle votazioni}
\label{art:2.3}

Le votazioni si svolgono esclusivamente in una delle seguenti sedi:

\begin{itemize}
\item in seduta, durante le riunioni in presenza della Lega e in particolare nel corso dell'Asta Iniziale;
\item nel canale di messaggistica ufficiale della Lega, purché sia garantita la tracciabilità dei voti.
\end{itemize}

La votazione è indetta dal Presidente ovvero su richiesta di almeno un quarto dei membri del Consiglio. Non è richiesta alcuna forma particolare di convocazione: è sufficiente la pubblicazione della proposta nella sede prescelta, con indicazione dell'oggetto.

Ciascun membro dispone di ventiquattro ore dalla pubblicazione della proposta per esprimere il proprio voto favorevole, contrario o di astensione. Decorso tale termine senza risposta, il membro non prende parte alla votazione e non è computato tra i membri aventi diritto. Nelle votazioni in seduta il voto è espresso immediatamente e chi è assente non prende parte alla votazione.

Qualora almeno tre quarti dei membri del Consiglio richiedano congiuntamente una votazione avente ad oggetto la condotta del Presidente, questi è obbligato a indirla entro ventiquattro ore. Decorso inutilmente tale termine, la votazione si intende validamente indetta dai richiedenti e presieduta dal membro da essi designato; in tale ipotesi il Presidente è in ogni caso escluso dal voto.

\section{Validità della votazione}
\label{art:2.4}

La votazione è valida se, scaduto il termine, hanno preso parte più della metà dei membri del Consiglio. In difetto la proposta si intende respinta e resta in vigore la disciplina precedente.

\section{Conflitto di interessi ed esclusione dal voto}
\label{art:2.5}

È escluso dalla singola votazione il membro che ne sia destinatario diretto, ovvero che ne ricavi un vantaggio patrimoniale o sportivo diretto, individuale e non condiviso con gli altri partecipanti.

L'esclusione è disposta dall'Amministratore di Lega, sentiti gli altri membri, che possono segnalare situazioni di conflitto prima dell'apertura della votazione. Essa può essere riesaminata con le modalità di cui all'articolo \ref{art:2.7}.

Qualora la votazione abbia ad oggetto la condotta o la posizione dell'Amministratore di Lega, l'esclusione è disposta dal membro che presiede la votazione ai sensi dell'articolo \ref{art:2.3}.

Il Presidente può disporre la sospensione temporanea dal voto del membro che adotti comportamenti lesivi, utilizzi linguaggio offensivo o ponga in essere ostruzionismo deliberato. Anche tale decisione può essere riesaminata ai sensi dell'articolo \ref{art:2.7}.

\section{Finestra di modifica e stabilità della stagione}
\label{art:2.6}

Le deliberazioni che incidono sulle regole o sulla struttura del gioco sono assunte esclusivamente nella Prima Fase dell'Asta Iniziale (vedi \ref{art:4.3.1}) e restano immutabili per l'intera stagione. Vi rientrano le modifiche al presente regolamento, l'adozione o la revoca dei fattori e delle funzioni opzionali, le quote e il montepremi, il numero delle squadre e la modalità di gioco.

Ogni deliberazione della medesima natura assunta in un momento diverso acquista efficacia a decorrere dalla stagione successiva, ancorché approvata con la maggioranza richiesta.

In deroga a quanto precede, con votazione straordinaria è deliberata con efficacia immediata sulla stagione in corso la modifica resa necessaria da:

\begin{itemize}
\item un evento sopravvenuto non disciplinato dal presente regolamento o non imputabile ai partecipanti;
\item una disposizione la cui applicazione impedisca, blocchi o renda impraticabile lo svolgimento del gioco.
\end{itemize}

Il mero disaccordo sul merito di una regola, o la sua sopravvenuta impopolarità, non integrano alcuna delle due ipotesi: in tali casi la modifica è deliberata ma acquista efficacia dalla stagione successiva. La delibera straordinaria è limitata a quanto strettamente necessario e, nell'ipotesi di cui al primo punto, decade con l'evento che l'ha determinata.

Non costituiscono modifica ai sensi del presente articolo le deliberazioni riguardanti le persone e la continuità della Lega (sanzioni, sospensioni, esclusioni, ammissioni, attribuzione e revoca delle cariche, nonché le misure di cui all'articolo \ref{art:1.5}), che restano deliberabili in qualsiasi momento, ancorché incidano sulla quota del singolo partecipante. Resta altresì ferma la disciplina temporale propria dell'articolo \ref{art:11.4}.

\section{Atti esecutivi e riesame}
\label{art:2.7}

Costituiscono atti esecutivi le decisioni che il presente regolamento rimette al Presidente o all'Amministratore di Lega, in quanto attuative di criteri già stabiliti. Esse sono adottate senza necessità di votazione e sono immediatamente efficaci.

Ogni atto esecutivo può essere riesaminato su richiesta anche di un solo membro, avanzata entro settantadue ore in presenza o a distanza: non è richiesta alcuna forma né alcun atto scritto, è sufficiente chiederne il riesame. Sulla richiesta decide il Consiglio. La richiesta non sospende l'atto, salvo diversa decisione del Consiglio.


\chapter{Sistema Sanzionatorio}
\label{cap:sistema-sanzionatorio}

\section{Principi generali}
\label{art:3.1}

Il presente titolo disciplina le conseguenze dei comportamenti contrari al regolamento, alla correttezza sportiva o allo spirito della competizione. Ogni sanzione deve essere proporzionata alla gravità del fatto e applicata allo stesso modo nei confronti di tutti.

\section{Infrazioni}
\label{art:3.2}

Costituisce infrazione ogni condotta contraria al presente regolamento o alla correttezza sportiva, e in particolare la mancata consegna reiterata della formazione, i comportamenti o i linguaggi offensivi, il tentativo di alterare l'esito delle competizioni e il mancato rispetto degli obblighi economici.

Per le infrazioni alle quali il regolamento ricollega una sanzione automatica si applica esclusivamente l'articolo \ref{art:3.4}.

\section{Sanzioni}
\label{art:3.3}

In funzione della gravità e della recidiva il Consiglio può disporre il richiamo ufficiale, la penalizzazione in crediti, la penalizzazione in classifica, la maggiorazione della quota di partecipazione fino al 50\% e, nei casi più gravi, l'esclusione definitiva dalla Lega.

In caso di esclusione la squadra prosegue la stagione ai sensi dell'articolo \ref{art:1.5} e i risultati già maturati restano validi. Qualora la quota stagionale non risulti versata e non vi siano modalità efficaci di riscossione, l'importo dovuto è ripartito in parti uguali tra gli altri membri.

\section{Procedura}
\label{art:3.4}

Le sanzioni sono deliberate dal Consiglio. L'interessato è informato preventivamente e dispone di ventiquattro ore per presentare le proprie difese; la votazione è indetta dopo la scadenza di tale termine e, in caso di silenzio, si procede comunque. L'esito è comunicato nel canale ufficiale della Lega e la sanzione è definitiva.

Le sanzioni espressamente predeterminate dal presente regolamento, quale quella prevista dall'articolo \ref{art:10.2}, si applicano di diritto al verificarsi del presupposto, senza necessità di deliberazione, e costituiscono l'unica conseguenza del fatto cui si riferiscono, salvo quanto diversamente previsto dall'articolo medesimo.


\chapter{Il Gioco}
\label{cap:il-gioco}

\section{Modalità di gioco}
\label{art:4.1}

A decorrere dalla stagione 2026/2027 il campionato si svolge secondo il sistema Mantra, disciplinato dal regolamento ufficiale di \emph{Fantacalcio.it} e applicato attraverso la piattaforma “Fantacalcio.it”. Punteggi, valutazioni e funzionalità tecniche seguono quanto previsto dal sito di riferimento, salvo le modifiche e le integrazioni del presente regolamento.

Il passaggio dal sistema Classic al sistema Mantra è deliberato dal Consiglio ed è soggetto alla finestra di modifica di cui all'articolo \ref{art:2.6}.

\section{Composizione della rosa}
\label{art:4.2}

Ogni squadra deve costituire una rosa di 30 calciatori, di cui almeno 3 portieri. Il sistema Mantra non contingenta gli altri ruoli: fermo quel minimo, l'assortimento è rimesso alle scelte tattiche di ciascun fantallenatore. La rosa si compone durante l'Asta Iniziale e va rispettata in ogni momento della stagione.

Ciascun calciatore può appartenere a una sola rosa e, una volta aggiudicato, non è più disponibile per le altre squadre fino al suo eventuale svincolo.

I dodici ruoli Mantra (Por, Dc, B, Dd, Ds, E, M, C, T, W, A, Pc) e le quotazioni sono assegnati da \emph{Fantacalcio.it}, che ne è fonte unica e non personalizzabile dalla Lega; molti calciatori ne ricoprono più di uno. Il listone è sottoposto a una sola revisione, alla chiusura del calciomercato estivo, e tale revisione può soltanto aggiungere ruoli, mai toglierne. Dopo di essa ruoli e quotazioni restano invariati per l'intera stagione.

\section{Asta iniziale}
\label{art:4.3}

\subsection{Prima fase: regole e votazioni}
\label{art:4.3.1}

Prima dell'avvio dell'asta il Consiglio si riunisce e delibera in seduta (vedi \ref{art:2.3}) sulle modifiche regolamentari e sulle scelte di stagione, unica sede in cui possono essere assunte con efficacia per la stagione in corso (vedi \ref{art:2.6}), nonché sul metodo di svolgimento dell'asta, valido anche per le aste successive.

Le scelte in vigore sono indicate in ciascuna regola del presente regolamento, che va aggiornato nell'articolo di volta in volta interessato ogniqualvolta il Consiglio ne deliberi la modifica.

\subsection{Seconda fase: composizione delle rose}
\label{art:4.3.2}

Ogni squadra dispone di un budget iniziale di 1.000 crediti e ogni calciatore parte da una base d'asta di 1 credito.

L'asta si articola in due fasi successive: prima i soli portieri, poi tutti gli altri calciatori senza distinzione di ruolo. Entro tale scansione il metodo di chiamata è deliberato dal Consiglio fra i seguenti:

\begin{itemize}
\item scorrimento alfabetico dei cognomi;
\item chiamata casuale, tramite il software “FantaAsta” ovvero mediante la «Ruota della Fortuna» (vedi \ref{art:1.6});
\item chiamata diretta libera dei fantallenatori, secondo l'ordine determinato ai sensi dell'articolo \ref{art:1.6}.
\end{itemize}

La Lega adotta la chiamata casuale tramite il software “FantaAsta”. Il metodo prescelto vale per ogni asta della stagione.

\subsection{Obbligo di completamento della rosa}
\label{art:4.3.3}

L'asta non si conclude finché ciascun partecipante non abbia raggiunto il numero esatto di 30 calciatori, con almeno 3 portieri. L'obbligo riguarda la rosa nel suo complesso e si verifica al termine dell'asta.

\subsection{Chiamata diretta integrativa}
\label{art:4.3.4}

Al termine dello scorrimento dell'elenco, in ciascuna delle due fasi, le squadre che non abbiano completato la rosa proseguono con una fase integrativa di chiamata diretta, scegliendo liberamente i calciatori mancanti secondo l'ordine determinato ai sensi dell'articolo \ref{art:1.6}.

\subsection{Correzione errori e svincoli}
\label{art:4.3.5}

Al termine di ciascuna delle due fasi i fantallenatori possono svincolare i calciatori non desiderati e recuperare per intero il credito speso. Lo svincolo va dichiarato subito ed è irrevocabile.

Il calciatore svincolato torna acquistabile solo tramite nuova asta, aperta anche agli altri fantallenatori; se nessuno rilancia, chi lo aveva svincolato può riprenderlo al prezzo originario.

\subsection{Errori di assegnazione}
\label{art:4.3.6}

L'aggiudicazione di un calciatore già appartenente a un'altra rosa, oppure assente dal listone ufficiale, è nulla: il credito speso è restituito per intero e il calciatore torna disponibile all'acquisto.

\section{Schieramento della formazione}
\label{art:4.4}

La formazione va schierata entro e non oltre i 15 minuti precedenti l'inizio della prima partita del turno di Serie A. In caso di guasto o malfunzionamento dell'applicazione va scritta entro il medesimo termine nel canale di messaggistica ufficiale della Lega, e l'Amministratore di Lega la carica nell'app Leghe appena il servizio torna disponibile; fa fede l'orario del messaggio.

Ciascuna formazione titolare comprende 11 calciatori, schierati secondo uno degli undici schemi previsti dal sistema Mantra: 3-4-3, 3-4-1-2, 3-4-2-1, 3-5-2, 3-5-1-1, 4-3-3, 4-3-1-2, 4-4-2, 4-1-4-1, 4-4-1-1, 4-2-3-1.

La griglia degli schemi, con i ruoli ammessi in ciascuna posizione, è quella pubblicata da \emph{Fantacalcio.it} e la Lega non può modificarla. Oltre al portiere, ogni schema impiega sempre 5 calciatori di movimento di stampo difensivo (Dd, Ds, Dc, B, E, M) e 5 di stampo offensivo (C, T, W, A, Pc).

La panchina è composta da 12 calciatori, di cui almeno un portiere, e va schierata in ordine di preferenza assoluta: entra per primo chi è inserito per primo, sempre che il cambio risulti tatticamente possibile. L'ordine della panchina governa l'intero meccanismo delle sostituzioni (vedi \ref{art:4.5}).

Schierare un calciatore in una posizione diversa dal proprio ruolo è consentito, ove la piattaforma lo permetta, con un malus di 1 punto per ciascun calciatore adattato; alcune collocazioni restano comunque inibite in fase di inserimento. Il ricorso allo schieramento fuori posizione è riservato alle emergenze di rosa.

In caso di mancata consegna della formazione, il sistema recupera l'ultima formazione valida disponibile.

\section{Sostituzioni}
\label{art:4.5}

Sono consentite fino a cinque sostituzioni per giornata. Il portiere è sempre il primo a essere sostituito quando il titolare risulta assente, e il subentro della riserva consuma una delle sostituzioni disponibili.

Per i calciatori di movimento le sostituzioni non avvengono una per volta ma in blocco: il sistema conta i titolari assenti o senza voto e sceglie quale combinazione di panchinari far subentrare. Le soluzioni possibili sono di tre tipi:

\begin{itemize}
\item ottimale, quando rispetta lo schema schierato e non richiede alcun calciatore fuori posizione;
\item efficiente, quando evita i malus ma impone un cambio modulo;
\item adattata, quando ricostruisce uno schema valido solo collocando fuori posizione uno o più calciatori, con il malus di 1 punto per ciascuno.
\end{itemize}

In tutte e tre le modalità di sostituzione le soluzioni adattate, essendo le uniche che comportano malus, sono l'ultima risorsa, e a parità di rango prevale sempre la combinazione comandata dai calciatori inseriti per primi in panchina. Cambia soltanto il peso dato allo schema di partenza:

\begin{itemize}
\item \textbf{BASIC}: prima le ottimali, poi le efficienti con cambio modulo, infine le adattate;
\item \textbf{EASY}: prima le ottimali, poi le adattate, sempre e solo nello schema scelto; il cambio modulo è escluso, e la formazione nasce e muore con il modulo iniziale;
\item \textbf{MASTER}: ottimali ed efficienti valgono allo stesso modo e comanda per intero l'ordine panchina, poi le adattate, anche qui senza alcuna priorità dello schema iniziale.
\end{itemize}

La Lega adotta la modalità MASTER.

Il calciatore che riceve un'ammonizione senza ottenere voto è trattato come qualsiasi altro calciatore senza voto e viene sostituito, in luogo dell'altra opzione offerta dalla piattaforma, che lo lascerebbe in campo con voto 5,5. La Lega non assegna voti d'ufficio a chi resta senza voto.

Qualora nessuna soluzione risulti praticabile, il sistema riprende il calcolo con un calciatore in meno e la squadra gioca in inferiorità numerica.

\section{Recidiva nella mancata consegna della formazione}
\label{art:4.6}

Qualora un partecipante ometta la consegna della formazione per più di tre giornate complessive, anche non consecutive, durante una stessa stagione, il Presidente potrà avviare una procedura sanzionatoria ai sensi del Titolo \ref{cap:sistema-sanzionatorio} del presente regolamento.


\chapter{Modalità di Calcolo}
\label{cap:modalita-calcolo}

\section{Redazione e fonte dei dati}
\label{art:5.1}

Fra le tre fonti di voto disponibili sulla piattaforma, ossia la redazione di \emph{Fantacalcio.it}, il Voto Statistico FG e la Redazione Italia, la Lega adotta la redazione di \emph{Fantacalcio.it}, che vale sia per le valutazioni in pagella sia per l'assegnazione di bonus e malus. Il cambio di fonte è deliberato nella Prima Fase dell'Asta Iniziale (vedi \ref{art:2.6}).

Ruoli, quotazioni e statistiche provengono in ogni caso da \emph{Fantacalcio.it}, che ne è fonte unica.

\section{Soglie gol}
\label{art:5.2}

La soglia per l'assegnazione del primo gol fantacalcistico è fissata a 66 punti. I gol successivi si ottengono al raggiungimento di incrementi di 4 punti rispetto alla soglia precedente. Di seguito un esempio esplicativo:

\begin{itemize}
    \item 66 punti → 1 gol
    \item 70 punti → 2 gol
    \item 74 punti → 3 gol
    \item 78 punti → 4 gol
    \item etc.
\end{itemize}

\section{Bonus e Malus}
\label{art:5.3}

I bonus e malus applicati ai calciatori, da sommare o sottrarre al voto in pagella, sono i seguenti:

\begin{itemize}
    \item +3 per ogni rete segnata
    \item +1 per ogni assist, compreso quello da fermo (vedi \ref{art:5.3.1})
    \item +3 per ogni rigore parato
    \item +1 per il portiere che mantiene la porta inviolata
    \item +1 per ogni gol vittoria
    \item +0,5 per ogni gol pareggio
    \item -2 per ogni autogol
    \item -1 per ogni rete subita
    \item -3 per ogni rigore sbagliato
    \item -0,5 per ogni ammonizione
    \item -1 per ogni espulsione
    \item -1 per ogni calciatore schierato fuori posizione (vedi \ref{art:4.4})
\end{itemize}

Nel sistema Mantra la tabella dei bonus e malus distingue esclusivamente tra portieri e calciatori di movimento.

\subsection{Assist}
\label{art:5.3.1}

Dalla stagione 2026/2027 l'assist è assegnato per ogni passaggio volontario che porta un compagno a segnare. L'unico criterio è la volontarietà del passaggio: l'entità del vantaggio procurato non decide più se l'assist spetti o meno, ma serve soltanto a classificarlo come soft, standard o gold.

La Lega non adotta i quality assist: le tre categorie valgono tutte +1, e altrettanto vale l'assist da fermo. Un assist è un assist, quale che ne sia la qualità.

I valori qui indicati sono modificabili con deliberazione assunta nella Prima Fase dell'Asta Iniziale (vedi \ref{art:2.6}).

\section{D-Factor}
\label{art:5.4}

Il D-Factor, per esteso Fattore Difensivo, premia il rendimento del pacchetto difensivo. La Lega lo adotta nella versione 5+1, comprensiva del portiere.

Il D-Factor si calcola come segue:

\begin{itemize}
\item si considerano i 5 calciatori di movimento di stampo difensivo presenti in campo (ruoli Dc, B, Dd, Ds, E, M); qualora essi siano più di cinque, si prendono i 5 con il voto migliore;
\item almeno 3 dei cinque devono avere ruolo Dc, B, Dd o Ds;
\item a essi si aggiunge il portiere, per un totale di 6 calciatori;
\item si calcola la media aritmetica dei sei voti in pagella, al netto di bonus e malus.
\end{itemize}

Il D-Factor non si attiva qualora non si raggiunga il numero di calciatori necessario al calcolo, circostanza che può verificarsi solo giocando in inferiorità numerica.

La media così ottenuta viene convertita secondo il seguente schema:

\begin{itemize}
    \item Media < 5.00: -3 punti
    \item 5.00 ≤ Media < 5.25: -2 punti
    \item 5.25 ≤ Media < 5.50: -1 punto
    \item 5.50 ≤ Media < 5.75: -0.5 punti
    \item 5.75 ≤ Media < 6.00: 0 punti
    \item 6.00 ≤ Media < 6.25: +0.5 punti
    \item 6.25 ≤ Media < 6.50: +1 punto
    \item 6.50 ≤ Media < 6.75: +2 punti
    \item 6.75 ≤ Media < 7.00: +3 punti
    \item 7.00 ≤ Media < 7.25: +4 punti
    \item 7.25 ≤ Media < 7.50: +5 punti
    \item Media ≥ 7.50: +6 punti
\end{itemize}

L'adozione del D-Factor e del criterio di conversione sopra indicato è deliberata nella Prima Fase dell'Asta Iniziale (vedi \ref{art:2.6}) e resta immutabile per l'intera stagione.

\section{Bonus Capitano}
\label{art:5.5}

Il bonus Capitano è un punteggio aggiuntivo calcolato sul voto in pagella, al netto di bonus e malus, del calciatore designato capitano. Ogni fantallenatore indica un capitano e un vice-capitano all'atto della compilazione della formazione.

La piattaforma prevede due varianti: capitano fisso per l'intera stagione oppure modificabile giornata per giornata. La Lega adotta la seconda.

È eleggibile qualsiasi calciatore schierato in formazione, indipendentemente dal ruolo Mantra ricoperto.

Qualora il capitano non scenda in campo ovvero non ottenga voto, il bonus è calcolato sul vice-capitano. Se neppure quest'ultimo ottiene voto, il bonus non è assegnato.

Il bonus Capitano è assegnato come segue:

\begin{itemize}
    \item Voto ≤ 4.5: -1 punto
    \item Voto = 5.0: -0.5 punti
    \item Voto = 5.5: 0 punti
    \item Voto = 6.0: +0.5 punti
    \item Voto = 6.5: +1 punto
    \item Voto = 7.0: +1.5 punti
    \item Voto ≥ 7.5: +2 punti
\end{itemize}

Il bonus Capitano è deliberato nella Prima Fase dell'Asta Iniziale (vedi \ref{art:2.6}) e resta immutabile per l'intera stagione.  
\textbf{Nota bene:} i bonus e i malus del giocatore designato capitano non vengono raddoppiati.

\section{Funzione “Switch”}
\label{art:5.6}

La funzione “Switch” consente di garantire la titolarità di un determinato calciatore in formazione. Se il giocatore selezionato non dovesse scendere in campo dal primo minuto nella partita reale, verrà automaticamente sostituito dal panchinaro preselezionato.

Si tratta di un'inversione condizionata tra un titolare e una riserva, attivata esclusivamente in caso di non titolarità effettiva.

Lo Switch non costituisce una sostituzione: opera sulla formazione schierata prima di qualsiasi processo sostitutivo e non consuma alcuna delle sostituzioni disponibili di cui all'articolo \ref{art:4.5}. Una volta operata l'inversione, il calcolo prosegue con le ordinarie sostituzioni, secondo l'ordine della panchina.

Esistono due varianti della funzione:

\begin{itemize}
    \item \textbf{Switch Basic:} consente la sostituzione solo tra calciatori che condividano almeno un ruolo Mantra;
    \item \textbf{Switch Plus:} consente anche la sostituzione tra ruoli diversi, a condizione che lo schema risultante sia valido secondo le regole di gioco.
\end{itemize}

La Lega adotta la funzione nella variante Basic. La funzione e la variante sono deliberate nella Prima Fase dell'Asta Iniziale (vedi \ref{art:2.6}) e restano immutabili per l'intera stagione.

\section{Partite sospese, rinviate o anticipate}
\label{art:5.7}

In caso di partite rinviate o anticipate di oltre 7 giorni rispetto alla data ufficialmente prevista per la giornata di campionato, oppure recuperate dopo l'inizio della giornata successiva, verrà assegnato un voto d'ufficio pari a 6 a tutti i giocatori coinvolti.

Per le partite sospese si procederà come segue:

\begin{itemize}
    \item Se la sospensione avviene nel primo tempo e il recupero avviene dopo l'inizio della giornata successiva, tutti i giocatori riceveranno voto d'ufficio pari a 6;
    \item Se la sospensione avviene nel primo tempo ma la partita è recuperata prima della giornata successiva, si attenderà il recupero della partita;
    \item Se la sospensione avviene nel secondo tempo si attenderà il recupero della partita.
\end{itemize}

In presenza di casi eccezionali (es. infortuni o squalifiche sopravvenuti), si seguiranno le indicazioni fornite da \emph{Fantacalcio.it}.

Se il numero totale di partite rinviate è pari o superiore a tre all'interno della stessa giornata, si attenderà il recupero delle stesse, senza ricorrere al voto d'ufficio.


\chapter{Mercato}
\label{cap:mercato}

\section{Prima asta di riparazione (opzionale)}
\label{art:6.1}

Qualora l'Asta Iniziale si svolga prima della chiusura del calciomercato estivo, al termine della sessione di mercato si tiene una prima asta di riparazione per aggiornare le rose agli ultimi trasferimenti.

Ogni squadra partecipa con i crediti residui dell'Asta Iniziale, sommati a quelli recuperati con gli svincoli dichiarati prima della sessione. L'asta si svolge nelle medesime due fasi previste dall'articolo \ref{art:4.3.2}, prima i portieri e poi tutti gli altri calciatori, e con il metodo di chiamata adottato dalla Lega.

Prima dell'inizio ciascun fantallenatore comunica, per l'intera rosa, i calciatori che intende svincolare: costoro diventano subito liberi e i crediti sono restituiti per intero. Il calciatore svincolato torna acquistabile solo tramite nuova asta e, in mancanza di rilanci, al prezzo originario.

Ogni squadra deve inoltre risultare in regola, prima dell'asta, con i criteri delle Priorità (vedi \ref{art:6.4}), ove applicabili.

\section{Seconda asta di riparazione (obbligatoria)}
\label{art:6.2}

La seconda asta di riparazione si tiene al termine del calciomercato invernale ed è obbligatoria per tutte le squadre. Valgono le regole dell'Asta Iniziale (vedi \ref{art:4.3}), salvo diversa deliberazione del Consiglio: si utilizzano i crediti residui della stagione, compresi quelli derivanti dagli svincoli, va comunicata in anticipo l'intera lista dei calciatori da svincolare e le rose vanno regolarizzate secondo le Priorità (vedi \ref{art:6.4}) e le sessioni di scambio (vedi \ref{art:6.3}) prima dell'inizio.

\section{Sessioni di scambi}
\label{art:6.3}

Durante fasi specifiche della stagione, il Presidente può indire sessioni ufficiali di scambio tra squadre quale atto esecutivo (vedi \ref{art:2.7}), previa comunicazione a tutti i partecipanti.

\textbf{Fasi abilitate agli scambi:}
\begin{itemize}
\item Periodi di sosta per le nazionali;
\item Periodi compresi tra un'asta e la giornata successiva.
\item Eventuali sessioni aggiuntive deliberate dal Consiglio di Lega.
\end{itemize}

\textbf{Regole generali:}
\begin{itemize}
\item Gli scambi sono liberi quanto ai ruoli: qualsiasi calciatore può essere scambiato con qualsiasi altro. Sono tuttavia richieste due condizioni: la parità numerica fra le due parti (non è ammesso, ad esempio, scambiare due calciatori con uno solo) e la permanenza in ciascuna rosa di almeno 3 portieri. È pertanto nullo lo scambio che porti una delle due squadre al di sotto di tale minimo, ai sensi dell'articolo \ref{art:4.2};
\item Il valore d'acquisto segue il calciatore trasferito;
\item È ammesso includere una somma di crediti come conguaglio tra le parti, purché venga dichiarata nel dettaglio.
\end{itemize}

Ogni tentativo di elusione o simulazione può comportare l'applicazione delle sanzioni di cui al Titolo \ref{cap:sistema-sanzionatorio}.

\textbf{Procedura obbligatoria:}\\
Uno dei due fantallenatori pubblica lo scambio nel canale ufficiale della Lega indicando i calciatori coinvolti, i crediti originari d'acquisto e l'eventuale conguaglio; l'altro conferma nella stessa conversazione. Il Presidente verifica e ufficializza quale atto esecutivo (vedi \ref{art:2.7}), con messaggio pubblico e aggiornamento dei crediti. Lo scambio è valido solo dopo l'ufficializzazione.

\textbf{Limiti tecnici:}\\
Nel caso in cui uno scambio comporti il superamento del limite massimo di 1.000 crediti totali, l'Amministratore di Lega riduce proporzionalmente il valore dei calciatori acquisiti dalla squadra eccedente, fino a riportare il budget entro il limite. L'operazione costituisce atto esecutivo (vedi \ref{art:2.7}).

\section{Le priorità}
\label{art:6.4}

Il sistema delle “Priorità” è applicabile in caso di uscita di un calciatore dalla Serie A. Il fantallenatore che ne detiene i diritti può usufruire di una sostituzione agevolata.

\textbf{Condizioni per l'attivazione della priorità:}
\begin{itemize}
\item Il calciatore deve risultare ufficialmente trasferito fuori dalla Serie A;
\item La sostituzione è ammessa con un calciatore svincolato di qualsiasi ruolo, purché la sua quotazione attuale sia pari o inferiore a quella del giocatore uscito, calcolata al momento dell'ufficialità della cessione;
\item Il giocatore scelto deve essere presente nella lista svincolati dalla data dell'ultima asta ufficiale svolta.
\end{itemize}

\textbf{Modalità:}

Il nuovo calciatore verrà acquistato al medesimo prezzo del giocatore sostituito; l'operazione deve essere comunicata e approvata dal Presidente, che ne verifica i criteri quale atto esecutivo (vedi \ref{art:2.7}).

\section{Mercato prestiti}
\label{art:6.5}

Dalla stagione 2026/2027 la piattaforma consente di cedere un calciatore per un periodo determinato, non inferiore a tre giornate e non superiore al termine della stagione, in cambio di uno o più calciatori ovvero di crediti da spendere nelle sessioni di mercato successive. Alla scadenza il calciatore torna in via automatica alla squadra di provenienza.

Il prestito è ammesso nelle stesse fasi previste per gli scambi e segue la medesima procedura di comunicazione e ufficializzazione (vedi \ref{art:6.3}). In deroga a quanto ivi previsto non è richiesta la parità numerica fra le parti, poiché il prestito può essere anche asimmetrico. Restano fermi il minimo di 3 portieri in ciascuna rosa e il limite dei 1.000 crediti complessivi.

La Lega adotta il mercato prestiti; l'eventuale revoca è deliberata nella Prima Fase dell'Asta Iniziale (vedi \ref{art:2.6}).


\chapter{Gestione Infortuni e Indisponibilità}
\label{cap:gestione-infortuni}

\section{Definizioni generali}
\label{art:7.1}

Ai fini del presente regolamento:

\begin{itemize}
\item Si definisce \textbf{infortunato} un calciatore che, a causa di una condizione fisica o mentale accertata, non risulta convocabile per un numero significativo di partite consecutive.
\item Si definisce \textbf{indisponibile} un calciatore escluso dalle convocazioni per ragioni non mediche, tra cui motivi tattici, disciplinari, societari o di altra natura (es. fuori rosa, sospeso, squalificato).
\end{itemize}

La gestione dei casi dipende dalla durata della non convocazione e dalla sua causa.

\section{Infortuni di durata inferiore ai 4 mesi}
\label{art:7.2}

Se un calciatore in rosa subisce un infortunio la cui durata effettiva (compresa tra la data del primo stop e quella della prima nuova convocazione) è inferiore a 4 mesi, il fantallenatore non potrà sostituire il giocatore né ricorrere ad alcun meccanismo di supporto.

\section{Infortuni di durata pari o superiore ai 4 mesi}
\label{art:7.3}

Nel caso in cui la durata dell'infortunio sia pari o superiore a 4 mesi, oppure venga ufficialmente dichiarata la fine anticipata della stagione per il calciatore in questione, il fantallenatore ha diritto a sostituirlo tramite la lista degli svincolati.

\textbf{Condizioni per la sostituzione:}
\begin{itemize}
\item È sufficiente che una diagnosi ufficiale, pubblicata da fonti attendibili (es. società calcistica, Fantacalcio.it, principali testate sportive), indichi un'assenza prevista pari o superiore a 4 mesi;
\item La sostituzione è approvata dal Presidente quale atto esecutivo (vedi \ref{art:2.7}); una volta approvata non sarà più possibile annullarla, anche in caso di rientro anticipato del calciatore;
\item Il calciatore sostitutivo può essere di qualsiasi ruolo e deve essere presente nella lista svincolati dalla data dell'ultima asta ufficiale;
\item La quotazione di riferimento del giocatore da sostituire è quella in vigore il giorno dell'infortunio;
\item Il sostituto deve avere una quotazione attuale pari o inferiore a tale valore ed è acquisito al medesimo prezzo del giocatore sostituito.
\end{itemize}

 \textbf{Limite temporale:}

Il diritto alla sostituzione può essere esercitato fino al 31 marzo compreso. Dal 1° aprile in poi, nessuna sostituzione sarà più ammessa, anche in presenza di infortuni gravi o diagnosi di stagione conclusa.

\section{Calciatori indisponibili}
\label{art:7.4}

I calciatori considerati indisponibili per ragioni non mediche (es. esclusione per scelta tecnica, fuori rosa, squalifica prolungata, sospensione contrattuale) non possono essere sostituiti.\\
Il fantallenatore dovrà attendere il reintegro del calciatore o provvedere durante una regolare finestra di mercato.

\section{Emergenze sanitarie e situazioni eccezionali}
\label{art:7.5}

Le indisponibilità determinate da malattie infettive, epidemie o pandemie (quale, a titolo di esempio principale, la positività al virus SARS-CoV-2, COVID-19), nonché da patologie ad oggi sconosciute o di nuova insorgenza, sono gestite come ordinari infortuni in funzione della durata effettiva dell'indisponibilità: si applicano pertanto tutte le disposizioni del presente Titolo.

Qualora l'evento sanitario assuma dimensioni tali da incidere sul regolare svolgimento del campionato di Serie A, prevalgono le disposizioni del Titolo \ref{cap:interruzione-campionato}.

Ove la situazione non risulti riconducibile ad alcuna delle ipotesi precedenti, il Presidente indice una votazione straordinaria (vedi \ref{art:2.6}) per deliberare le misure necessarie.


\chapter{Interruzione del Campionato di Serie A}
\label{cap:interruzione-campionato}

\section{Sospensione straordinaria del campionato}
\label{art:8.1}

In caso di sospensione imprevista o straordinaria del campionato di Serie A (es. emergenze sanitarie, eventi eccezionali, motivazioni istituzionali), la Lega FantaScarrupat attenderà l'eventuale ripresa ufficiale delle partite per proseguire regolarmente la stagione.

Se la ripresa avviene con modalità differenti da quelle originarie (es. campionato compresso, tornei a eliminazione, playoff/playout), il gioco si interromperà all'ultima giornata completa disputata secondo il format iniziale, salvo diverse disposizioni del Consiglio di Lega assunte con votazione straordinaria (vedi \ref{art:2.6}).

In caso di ripresa posticipata, tutti i dati regolamentari (punteggi, penalità, scambi, sostituzioni, fattori, ecc.) restano validi, salvo diversa decisione del Consiglio di Lega assunta con votazione straordinaria (vedi \ref{art:2.6}).

\section{Validità delle competizioni e gestione delle quote}
\label{art:8.2}

Qualora il campionato non dovesse riprendere e venisse dichiarato concluso anticipatamente, si applicheranno le seguenti disposizioni, in base al numero di giornate effettivamente disputate; le frazioni di giornata si arrotondano all'unità superiore:

\textbf{a) Se la stagione raggiunge almeno il 60\% delle giornate previste:}
\begin{itemize}
\item La stagione è considerata valida;
\item I premi previsti verranno ricalcolati proporzionalmente in base al numero effettivo di giornate giocate;
\item Le quote eccedenti (cioè non utilizzate per premi e spese) saranno restituite oppure compensate sulla quota della stagione successiva;
\item L'albo d'oro viene regolarmente aggiornato, ad eccezione delle competizioni la cui finale unica non sia stata disputata, che non vengono conteggiate.
\end{itemize}

\textbf{b) Se la stagione non raggiunge il 60\% delle giornate previste:}
\begin{itemize}
\item La stagione è considerata annullata ai fini sportivi;
\item I premi vengono completamente annullati e non saranno distribuiti;
\item Le quote versate saranno integralmente restituite oppure compensate sulla stagione successiva;
\item L'albo d'oro non viene aggiornato: la stagione annullata non assegna alcun titolo.
\end{itemize}

\section{Sospensioni programmate e pause da calendario}
\label{art:8.3}
 
Le presenti disposizioni non si applicano nei casi in cui il campionato venga interrotto per motivi già previsti dal calendario ufficiale, come:
\begin{itemize}
\item Mondiali o Europei per nazionali;
\item Competizioni continentali (es. Coppa d'Africa, Copa América);
\item Olimpiadi o altri eventi internazionali ufficiali.
\end{itemize}

In questi casi, il gioco si sospende automaticamente in attesa della regolare ripresa del campionato, senza modifiche al regolamento o al calendario delle competizioni.

\section{Votazione Straordinaria}
\label{art:8.4}

In presenza di interruzioni impreviste, modifiche al format della Serie A o contesti non disciplinati da questo articolo, il Presidente ha facoltà di indire una votazione straordinaria (vedi \ref{art:2.6}) per deliberare su:
\begin{itemize}
\item Eventuale chiusura anticipata della stagione in corso;
\item Ricalcolo dei premi e delle quote in modalità alternativa;
\item Possibile inserimento o modifica di competizioni extra.
\end{itemize}


\chapter{Competizioni}
\label{cap:competizioni}

\section{Serie A}
\label{art:9.1}

Alla competizione “Serie A” partecipano tutte le 8 squadre della Lega.\\
Il torneo si svolge con girone asimmetrico e calendario generato in modo casuale, dal presidente di lega in presenza di almeno un altro membro del consiglio come testimone, tramite l'apposito software di Fantacalcio.it, durante la fase d'Asta Iniziale.

\textbf{Punteggio:}
\begin{itemize}
\item 3 punti per la vittoria;
\item 1 punto per il pareggio;
\item 0 punti per la sconfitta.
\end{itemize}

\textbf{Classifica finale e premi:}\\
Le prime 3 classificate sono considerate vincitrici e accedono alla distribuzione dei premi stagionali.

In caso di parità in classifica, si applicano i seguenti criteri nell'ordine:
\begin{itemize}
\item Somma dei fantapunti totalizzati in tutte le giornate;
\item Classifica avulsa;
\item Differenza reti;
\item Gol fatti;
\item Gol subiti.
\end{itemize}

\section{Champions League}
\label{art:9.2}
Partecipano tutte le 8 squadre, suddivise in due fasi:

\subsection{Fase a gironi}
\label{art:9.2.1}

Le squadre vengono divise in 2 gruppi da 4, generato in modo casuale insieme al calendario, dal presidente di lega in presenza di almeno un altro membro del consiglio come testimone, tramite l'apposito software di Fantacalcio.it, durante la fase d'Asta Iniziale.\\
Ogni squadra disputa 6 incontri, affrontando le avversarie del proprio gruppo.

\textbf{Punteggio gironi:}
\begin{itemize}
\item 3 punti per la vittoria;
\item 1 punto per il pareggio;
\item 0 punti per la sconfitta.
\end{itemize}

\textbf{Criteri in caso di parità nei gironi:}
\begin{itemize}
\item Somma dei fantapunti totalizzati nel girone;
\item Classifica avulsa;
\item Differenza reti;
\item Gol fatti;
\item Gol subiti.
\end{itemize}

Le prime due classificate di ciascun girone accedono alla fase successiva.

\subsection{Eliminazione diretta}
\label{art:9.2.2}

Le prime classificate sfidano le seconde del girone opposto con gare di andata e ritorno. La finale si gioca in partita secca.

\textbf{Criteri per determinare la vincente:}
\begin{itemize}
\item Somma dei punteggi delle due gare;
\item Supplementari (vedi \ref{art:9.5.1});
\item Rigori (vedi \ref{art:9.5.2});
\item Somma dei fantapunti totalizzati in Serie A;
\item Gol in trasferta.
\end{itemize}

\section{NBA}
\label{art:9.3}

Questa competizione coinvolge tutte le 8 squadre ed è suddivisa in:

\subsection{Regular Season}
\label{art:9.3.1}

Si disputa con girone asimmetrico di 14 giornate, con calendario generato in modo casuale, dal presidente di lega in presenza di almeno un altro membro del consiglio come testimone, tramite l'apposito software di Fantacalcio.it, durante la fase d'Asta Iniziale.

\textbf{Punteggio:}
\begin{itemize}
\item 3 punti per la vittoria;
\item 1 punto per il pareggio;
\item 0 punti per la sconfitta.
\end{itemize}

\textbf{Criteri di parità in classifica:}
\begin{itemize}
\item Somma dei fantapunti totalizzati nella Regular Season;
\item Classifica avulsa;
\item Differenza reti;
\item Gol fatti;
\item Gol subiti.
\end{itemize}

Tutte le 8 squadre si qualificano alla fase successiva.

\subsection{Play-Off}
\label{art:9.3.2}

Competizione a eliminazione diretta, con andata e ritorno fino alla semifinale e finale secca.

\textbf{Tabellone (basato sulla Regular Season):}
\begin{itemize}
    \item Match 1: 1ª vs 8ª
    \item Match 2: 4ª vs 5ª
    \item Match 3: 3ª vs 6ª
    \item Match 4: 2ª vs 7ª
\end{itemize}

\textbf{Semifinali:} 
\begin{itemize}
    \item Match 5: vincente Match 1 vs vincente Match 2
    \item Match 6: vincente Match 3 vs vincente Match 4
\end{itemize}
\textbf{Finale:} vincente delle due semifinali

\textbf{Criteri di vittoria:}
\begin{itemize}
\item Risultato totale tra andata e ritorno;
\item Supplementari (vedi \ref{art:9.5.1});
\item Rigori (vedi \ref{art:9.5.2});
\item Somma dei fantapunti totalizzati in Serie A;
\item Gol in trasferta.
\end{itemize}

\section{Supercoppa Italiana}
\label{art:9.4}
Questa competizione coinvolge un totale di quattro squadre tra quelle appartenenti alla Lega. Le partecipanti vengono determinate come segue:

\begin{itemize}
    \item Vincitore della competizione “Serie A” della stagione precedente;
    \item Vincitore della competizione “Champions League” della stagione precedente;
    \item Vincitore della competizione “NBA” della stagione precedente;
    \item Secondo classificato della competizione “Serie A” della stagione precedente.
\end{itemize}

Qualora una medesima squadra risulti idonea in base a più criteri, ovvero un posto resti comunque vacante, si procede nell'ordine:

\begin{itemize}
    \item il posto reso vacante è assegnato al finalista perdente della competizione che lo ha generato;
    \item qualora anche tale finalista risulti già qualificato, il posto è assegnato secondo l'ordine decrescente della classifica di “Serie A” della stagione precedente (3ª, 4ª, 5ª e così via), con esclusione delle squadre già qualificate.
\end{itemize}

La competizione si svolge secondo il formato ad eliminazione diretta, con partite di andata e ritorno nelle semifinali, e finale in gara secca.

\textbf{Tabellone delle semifinali (basato sui risultati della stagione precedente):}
\begin{itemize}
    \item Semifinale 1: Vincitore “Serie A” vs Vincitore “NBA”;
    \item Semifinale 2: Secondo classificato “Serie A” vs Vincitore “Champions League”.
\end{itemize}

Il tabellone si compone in base alla posizione occupata nell'elenco delle partecipanti, indipendentemente dal titolo che ha determinato la qualificazione di ciascuna squadra.

Il diritto a partecipare spetta al fantallenatore che lo ha conquistato e non alla squadra. Qualora egli lasci la Lega, il posto non passa a chi ne assume o ne eredita la squadra: esso si considera vacante ed è assegnato secondo i criteri di sostituzione sopra indicati, scorrendo l'ordine di accesso alla competizione.

\textbf{Criteri di vittoria:}
\begin{itemize}
\item Risultato totale tra andata e ritorno;
\item Supplementari (vedi \ref{art:9.5.1});
\item Rigori (vedi \ref{art:9.5.2});
\item Somma dei fantapunti totalizzati in Serie A;
\item Gol in trasferta.
\end{itemize}

\section{Supplementari e Rigori}
\label{art:9.5}

Le regole sui supplementari e i rigori si applicano esclusivamente alle fasi a eliminazione diretta delle competizioni in cui sono espressamente previste, ovvero Champions League, NBA (Play-Off) e Supercoppa Italiana.\\
Non sono previsti per le gare della fase a gironi, della Regular Season o per eventuali recuperi o spareggi non ufficiali.

\subsection{Supplementari}
\label{art:9.5.1}

Ogni squadra calcola un punteggio supplementare basato sulla media delle fantamedie complessive (comprensive di bonus e malus) dei migliori quattro panchinari, esclusi i portieri, non subentrati in campo.\\
Se meno di quattro giocatori ottengono voto valido, si integrano con voti d'ufficio pari a 5,5 fino a raggiungere il numero richiesto.

\begin{longtable}{|c|c|}
\hline
\textbf{Media complessiva} & \textbf{Gol supplementari} \\
\hline
0.00 - 6.49 & 0 gol \\ \hline
6.50 - 6.99 & 1 gol \\ \hline
7.00 - 7.49 & 2 gol \\ \hline
+0.50 ulteriore & +1 gol aggiuntivo \\
\hline

\caption{Gol x Media Supplementari}
\end{longtable}

\subsection{Rigori}
\label{art:9.5.2}

Se i supplementari non determinano una vincitrice, si procede con i rigori.\\
Il risultato di ciascun rigore è calcolato esclusivamente in base al voto puro (senza bonus o malus) del calciatore.

\begin{itemize}
\item Il calciatore segna se il voto puro è maggiore o uguale a 6;
\item Sbaglia se il voto è inferiore a 6.
\end{itemize}

Ogni squadra esegue 5 rigori. In caso di ulteriore parità si procede a oltranza, un rigore alla volta per squadra, seguendo l'ordine indicato e fino a esaurimento degli undici calciatori che compongono la formazione valida. Qualora, esauriti i tiratori disponibili, la parità persista, si applicano i successivi criteri previsti per la competizione di riferimento.

\textbf{Ordine di esecuzione dei rigori:}
\begin{itemize}
\item Punte centrali e attaccanti di raccordo (Pc, A) con voto valido, in ordine di schieramento;
\item Ali e trequartisti (W, T) con voto valido;
\item Centrocampisti (C, M) ed esterni difensivi (E) con voto valido;
\item Difensori (Dc, B, Dd, Ds) con voto valido;
\item Portiere con voto valido;
\item Riserve subentrate con voto valido, in ordine panchina.
\end{itemize}

Il calciatore polivalente è collocato nell'ordine secondo il ruolo effettivamente ricoperto nello schieramento.


\chapter{Quote, Penalità e Premi}
\label{cap:quote-penalita-premi}

\section{Quote di partecipazione}
\label{art:10.1}

Per prendere parte alla stagione, ogni squadra corrisponde una quota complessiva di 160 euro, versata al Tesoriere della Lega.\\
L'importo è frazionato in due rate uguali da 80 euro: la prima va pagata entro il giorno in cui si tiene la Seconda Asta di Riparazione, mentre la seconda è dovuta in occasione dell'ultima giornata di Serie A, quando i partecipanti sono soliti ritrovarsi per il consueto aperitivo di fine campionato, offerto dalle squadre vincitrici.\\
Al momento di ciascun pagamento il Tesoriere rilascia una semplice ricevuta, non fiscale ma nominativa, che costituisce l'unica prova ufficiale dell'avvenuto versamento.

\section{Penalità per ritardato pagamento}
\label{art:10.2}

La mancata corresponsione di una rata entro la relativa scadenza comporta, a carico della squadra inadempiente, una penalizzazione in classifica pari al numero di giornate di Serie A già disputate alla data di scadenza non rispettata.

La penalizzazione di cui al presente articolo costituisce l'unica conseguenza del ritardato pagamento. Il ricorso agli strumenti del Titolo \ref{cap:sistema-sanzionatorio} è ammesso esclusivamente in caso di mancato pagamento protratto oltre trenta giorni dalla scadenza; in tale ipotesi la penalizzazione già applicata è assorbita dalla sanzione deliberata.

\section{Premi}
\label{art:10.3}

Le quote annuali (8 squadre × € 160 = € 1 280) costituiscono il montepremi da distribuire come segue:

\begin{longtable}{|l|l|r|}
\hline
\textbf{Competizione} & \textbf{Posizione / Titolo} & \textbf{Premio (€)} \\
\hline
Serie A & 1ª & 350 \\
        & 2ª & 270 \\
        & 3ª & 200 \\ \hline
Champions League & Vincitrice & 190 \\ \hline
NBA & Vincitrice & 190 \\ \hline
Supercoppa Italiana & Vincitrice & 80 \\ \hline
\caption{Premi}
\end{longtable}

\section{Adeguamenti in caso di stagione abbreviata}
\label{art:10.4}

Le disposizioni dell'articolo \ref{art:8.2} si applicano anche al montepremi:

\begin{itemize}
\item \textbf{Stagione valida} (almeno il 60\% delle giornate giocate): i premi sono proporzionati alle giornate disputate; eventuali quote eccedenti sono restituite o compensate sulla quota della stagione successiva.
\item \textbf{Stagione annullata} (meno del 60\% delle giornate giocate): tutti i premi sono annullati; le quote versate sono restituite o compensate sulla stagione seguente.
\end{itemize}

\section{Fondo di riserva}
\label{art:10.5}

Eventuali residui non distribuiti (arrotondamenti, penalità pecuniarie rientrate, ecc.) confluiscono in un fondo di riserva gestito dal Tesoriere.\\
L'impiego del fondo richiede deliberazione del Consiglio e può essere destinato a:
\begin{itemize}
    \item spese straordinarie di Lega;
    \item incremento dei premi della stagione successiva;
    \item altro.
\end{itemize}

\section{Trofei delle competizioni}
\label{art:10.6}

Per ciascuna stagione la Lega provvede alla realizzazione dei trofei destinati alle competizioni di cui al Titolo \ref{cap:competizioni}, prodotti mediante stampa tridimensionale.

In occasione della Seconda Asta di Riparazione (vedi \ref{art:6.2}) il Tesoriere illustra al Consiglio i materiali necessari alla realizzazione e il costo complessivo che ne deriva.

Tale costo è ripartito in parti uguali tra tutti i partecipanti ed è dovuto in aggiunta alla quota di partecipazione di cui all'articolo \ref{art:10.1}, dalla quale resta distinto: esso non concorre alla formazione del montepremi (vedi \ref{art:10.3}) né incide sulla sua ripartizione.

Qualora il fondo di riserva (vedi \ref{art:10.5}) presenti una disponibilità, il Consiglio può deliberare di attingervi in tutto o in parte, riducendo in misura corrispondente la somma a carico dei partecipanti.

La determinazione dell'importo non costituisce modifica delle quote ai sensi dell'articolo \ref{art:2.6}, trattandosi della mera constatazione del costo dei materiali per la stagione in corso.

I trofei restano definitivamente di proprietà delle squadre che li hanno conquistati e non sono rimessi in palio nelle stagioni successive.


\chapter{Sito Web e Identità delle Fantasquadre}
\label{cap:sito-identita}

\section{Portale ufficiale}
\label{art:11.1}

Dal campionato 2022-2023 la Lega FantaScarrupat dispone di un portale istituzionale raggiungibile all'indirizzo \texttt{www.fantascarrupat.altervista.org}.\\
Il sito costituisce il centro informativo dell'intera Lega: ospita regolamento, calendari, rose, risultati e comunicati ufficiali, custodisce l'albo d'oro e l'archivio storico delle stagioni concluse e raccoglie contenuti multimediali e statistiche pubblicati dall'Amministratore di Lega quali atti esecutivi (vedi \ref{art:2.7}).\\
Il portale comprende inoltre un'area riservata, accessibile ai soli utenti abilitati, nella quale sono pubblicati i documenti e i materiali destinati ai membri del Consiglio di Lega.

\section{Gestione e responsabilità}
\label{art:11.2}

La gestione tecnica del portale spetta all'Amministratore di Lega.\\
La correttezza dei contenuti è invece responsabilità solidale di tutti i membri: ciascun membro, qualora rilevi inesattezze o lacune, è tenuto a segnalarle tempestivamente all'Amministratore di Lega e a proporre eventuali integrazioni.

\section{Nomi e loghi delle Fantasquadre}
\label{art:11.3}

Ogni nuovo Fantallenatore sceglie un nome e realizza un logo originali per la propria squadra; marchi, stemmi o elementi grafici già esistenti non possono essere riutilizzati, in ossequio alla normativa sul diritto d'autore.\\
Nome e logo vengono pubblicati sul portale e restano vincolanti finché non sia stata autorizzata una modifica secondo le disposizioni dell'articolo \ref{art:11.4}.

\section{Procedura di modifica}
\label{art:11.4}

Il Fantallenatore che intenda cambiare nome o logo compila l'apposito modulo “Richiesta cambi nome e logo” e lo inoltra nella chat ufficiale di Lega; l'Amministratore di Lega verifica la conformità ai requisiti e apre una votazione del Consiglio, che delibera con la maggioranza di cui all'articolo \ref{art:2.2}.\\
Le modifiche per motivi urgenti possono essere presentate in qualunque momento; quelle esclusivamente estetiche sono ammesse solo all'inizio della stagione e, se accolte, non potranno essere richieste di nuovo prima che siano trascorse due stagioni complete.\\
In caso di approvazione, l'Amministratore di Lega aggiorna il sito e il materiale ufficiale entro settantadue ore.

\section{Sanzioni}
\label{art:11.5}

Chi modifica nome o logo senza autorizzazione è obbligato a ripristinare immediatamente l'identità approvata.\\
In caso di recidiva si applica un aumento di cinque euro sulla quota di partecipazione, fatto salvo il potere del Consiglio di adottare pene più severe qualora la violazione persista.

\section{Tutela dei dati}
\label{art:11.6}

Il trattamento dei dati personali dei partecipanti è disciplinato dall'informativa privacy pubblicata sul portale all'indirizzo \texttt{www.fantascarrupat.altervista.org/public/privacy.php}, redatta ai sensi del Regolamento (UE) 2016/679 (GDPR) e del D.lgs. 196/2003 e successive modifiche. Titolare del trattamento è il soggetto ivi indicato.

I dati sono trattati esclusivamente per le finalità sportive e associative della Lega e non sono ceduti a terzi. La pubblicazione sul portale di nome, immagine e risultati di ciascun partecipante avviene previo consenso dell'interessato, revocabile in qualsiasi momento con le modalità indicate nell'informativa; la revoca comporta la rimozione dei dati dalle pagine pubbliche e non pregiudica la partecipazione alla Lega.

Ogni modifica sostanziale dell'informativa è comunicata al Consiglio di Lega.


\chapter{Firme e accettazione del regolamento}
\label{cap:firme}

Il regolamento è la fonte normativa primaria della Lega FantaScarrupat: tutte le decisioni tecniche e disciplinari devono conformarsi ad esso e, in caso di dubbio interpretativo, il suo testo prevale su prassi o consuetudini. Con la sottoscrizione della presente edizione, riferita alla stagione sportiva 2026/2027, i partecipanti dichiarano di averlo letto e compreso, di approvarlo senza riserve e di impegnarsi a rispettarlo.

La Lega FantaScarrupat è un'associazione amatoriale priva di scopo di lucro e il presente regolamento produce effetti esclusivamente interni tra i sottoscrittori. Le obbligazioni economiche in esso previste (quote, penalità e premi) attengono a un rapporto di gioco e, ai sensi dell'articolo 1933 del Codice Civile, non sono azionabili in giudizio: la loro osservanza è affidata alla lealtà dei partecipanti ed è presidiata unicamente dagli strumenti disciplinari di cui al Titolo \ref{cap:sistema-sanzionatorio}. Il documento non costituisce contratto associativo o societario, né rapporto di scommessa, e non produce effetti nei confronti di terzi.

\begin{table}[!ht]
\centering
\begin{tabular}{|c|p{4.5cm}|p{4.5cm}|p{2.8cm}|}
\hline
\textbf{N.} & \textbf{Nome e cognome} & \textbf{Firma} & \textbf{Data} \\
\hline
1 & Antonio Di Giorgio (P) & & \\
\hline
2 & Mario Castaldi & & \\
\hline
3 & Cristian Cecere & & \\
\hline
4 & Luigi Lucariello & & \\
\hline
5 & Andrea Lucariello & & \\
\hline
6 & Pasquale Lupoli & & \\
\hline
7 & Francesco Parolisi & & \\
\hline
8 & Lorenzo Parolisi & & \\
\hline
\end{tabular}
\caption{Firme}
\end{table}
